\documentclass[10pt,a4paper]{amsart}
\usepackage{amsmath,amssymb,amsthm,graphicx}
\usepackage{verbatim}
\newenvironment{reference}{\comment}{\endcomment}
%\newenvironment{reference}{}{}
\newenvironment{slogan}{\comment}{\endcomment}
\newenvironment{history}{\comment}{\endcomment}

% For commutative diagrams we use Xy-pic
\usepackage[all]{xy}

% We use 2cell for 2-commutative diagrams.
\xyoption{2cell}
% Theorem environments.
%
\theoremstyle{plain}
\newtheorem{theorem}[subsection]{Theorem}
\newtheorem{proposition}[subsection]{Proposition}
\newtheorem{lemma}[subsection]{Lemma}

\theoremstyle{definition}
\newtheorem{definition}[subsection]{Definition}
\newtheorem{example}[subsection]{Example}
\newtheorem{exercise}[subsection]{Exercise}
\newtheorem{situation}[subsection]{Situation}

\theoremstyle{remark}
\newtheorem{remark}[subsection]{Remark}
\newtheorem{remarks}[subsection]{Remarks}

\numberwithin{equation}{subsection}

% Macros
%
\def\lim{\mathop{\mathrm{lim}}\nolimits}
\def\colim{\mathop{\mathrm{colim}}\nolimits}
\def\Spec{\mathop{\mathrm{Spec}}}
\def\Hom{\mathop{\mathrm{Hom}}\nolimits}
\def\Ext{\mathop{\mathrm{Ext}}\nolimits}
\def\SheafHom{\mathop{\mathcal{H}\!\mathit{om}}\nolimits}
\def\SheafExt{\mathop{\mathcal{E}\!\mathit{xt}}\nolimits}
\def\Sch{\mathit{Sch}}
\def\Mor{\mathop{\mathrm{Mor}}\nolimits}
\def\Ob{\mathop{\mathrm{Ob}}\nolimits}
\def\Sh{\mathop{\mathit{Sh}}\nolimits}
\def\NL{\mathop{N\!L}\nolimits}
\def\CH{\mathop{\mathrm{CH}}\nolimits}
\def\proetale{{pro\text{-}\acute{e}tale}}
\def\etale{{\acute{e}tale}}
\def\QCoh{\mathit{QCoh}}
\def\Ker{\mathop{\mathrm{Ker}}}
\def\Im{\mathop{\mathrm{Im}}}
\def\Coker{\mathop{\mathrm{Coker}}}
\def\Coim{\mathop{\mathrm{Coim}}}

% Boxtimes
%
\DeclareMathSymbol{\boxtimes}{\mathbin}{AMSa}{"02}

%
% Macros for moduli stacks/spaces
%
\def\QCohstack{\mathcal{QC}\!\mathit{oh}}
\def\Cohstack{\mathcal{C}\!\mathit{oh}}
\def\Spacesstack{\mathcal{S}\!\mathit{paces}}
\def\Quotfunctor{\mathrm{Quot}}
\def\Hilbfunctor{\mathrm{Hilb}}
\def\Curvesstack{\mathcal{C}\!\mathit{urves}}
\def\Polarizedstack{\mathcal{P}\!\mathit{olarized}}
\def\Complexesstack{\mathcal{C}\!\mathit{omplexes}}
% \Pic is the operator that assigns to X its picard group, usage \Pic(X)
% \Picardstack_{X/B} denotes the Picard stack of X over B
% \Picardfunctor_{X/B} denotes the Picard functor of X over B
\def\Pic{\mathop{\mathrm{Pic}}\nolimits}
\def\Picardstack{\mathcal{P}\!\mathit{ic}}
\def\Picardfunctor{\mathrm{Pic}}
\def\Deformationcategory{\mathcal{D}\!\mathit{ef}}
\usepackage{cleveref}
\crefname{originallemma}{original lemma}{original lemmas}
\Crefname{originallemma}{Original lemma}{Original lemmas}
\makeatletter
\newlabel{lemma-self-centralizing-subfield}{{074V}{source}{}{sourcepointer-1}{}}
\newlabel{lemma-self-centralizing-subfield@cref}{{[originallemma][0][]074V}{[1][0][]source}{}{}{}}
\makeatother
\crefname{originallemma}{original lemma}{original lemmas}
\Crefname{originallemma}{Original lemma}{Original lemmas}
\makeatletter
\newlabel{lemma-when-tensor-is-equal}{{074U}{source}{}{sourcepointer-2}{}}
\newlabel{lemma-when-tensor-is-equal@cref}{{[originallemma][0][]074U}{[1][0][]source}{}{}{}}
\makeatother
\crefname{originaltheorem}{original theorem}{original theorems}
\Crefname{originaltheorem}{Original theorem}{Original theorems}
\makeatletter
\newlabel{theorem-centralizer}{{074T}{source}{}{sourcepointer-3}{}}
\newlabel{theorem-centralizer@cref}{{[originaltheorem][0][]074T}{[1][0][]source}{}{}{}}
\makeatother
\title{Splitting fields of a central simple algebra and similar algebras of the same degree}
\author{The Stacks Project Authors}
\date{stacks-project:brauer.tex:local-snapshot-20261012; GNU FDL 1.2 or later}
\begin{document}
\maketitle
\noindent\textit{Editorial note.} All delivered mathematical text is the author's own, copied from the pinned Stacks Project chapter file sites-cohomology.tex; no statement or proof step was rewritten, reordered or removed.\par\medskip
\noindent\textit{Editorial note.} The proof is the author's complete argument for both directions; every identification is either displayed or justified by a named preceding result, and the closing sentence 'so we win anyway' follows the displayed dimension computation.\par\medskip
\noindent\textit{Editorial note.} Every labelled reference occurring inside the counted spans is carried in assembly.external\_references as a non-load-bearing navigation pointer (same-chapter ones at their definition site, cross-chapter ones at the invocation inside the counted span), so no load-bearing mathematics is replaced by a stub. No bibliography citation occurs inside the counted spans.\par\medskip
\noindent\textit{Editorial note.} The same-name check for this card scanned every staged Stacks chapter of this repository -- all sources/stacks-*/pinned-original-source.tex files, the card's own chapter excluded -- and the labels found are recorded in dedup.same\_name\_cross\_chapter on the card.\par\medskip
\noindent\textit{Editorial note.} The preamble is the chapter's own retained minimal preamble (sources/stacks-sites-cohomology/retained-preamble.tex), a verbatim excerpt of the native preamble: the Xy-pic setup and the native theorem declarations with the mathematical macros.\par\medskip
\noindent\textit{Allowed standard prerequisites (editorial).}\begin{itemize}
\item Self centralizing subfields. Same chapter, Lemma lemma-self-centralizing-subfield. Supplies that the image of the field in the opposite algebra is its own centralizer, which identifies the tensor product with the centralizer of the ground field inside an endomorphism algebra.
\item The commutant and the Brauer group. Same chapter, Lemma lemma-when-tensor-is-equal. Supplies that the class of the commutant is the negative of the class of the algebra in the Brauer group, which gives similarity of the two algebras.
\item The dimension formula for the centralizer. Same chapter, Theorem theorem-centralizer. Supplies the identity relating the degrees of the algebra and of its commutant, from which the degree of the chosen algebra is computed.
\end{itemize}\medskip
\ifdefined\hypertarget\hypertarget{sourcepointer-1}{}\fi
\noindent\textit{Source reference.} Original lemma 074V: sources/stacks-brauer/pinned-original-source.tex, source lines 800--800. Navigation pointer invoked at proof line 828; the counted span contains the pointer only and the cited result is not reproduced inside it, so the pointer is registered as a non-load-bearing navigation stub in the established form for Stacks material. No bibliography citation is involved.\par\medskip
\ifdefined\hypertarget\hypertarget{sourcepointer-2}{}\fi
\noindent\textit{Source reference.} Original lemma 074U: sources/stacks-brauer/pinned-original-source.tex, source lines 784--784. Navigation pointer invoked at proof line 878; the counted span contains the pointer only and the cited result is not reproduced inside it, so the pointer is registered as a non-load-bearing navigation stub in the established form for Stacks material. No bibliography citation is involved.\par\medskip
\ifdefined\hypertarget\hypertarget{sourcepointer-3}{}\fi
\noindent\textit{Source reference.} Original theorem 074T: sources/stacks-brauer/pinned-original-source.tex, source lines 746--746. Navigation pointer invoked at proof line 793; the counted span contains the pointer only and the cited result is not reproduced inside it, so the pointer is registered as a non-load-bearing navigation stub in the established form for Stacks material. No bibliography citation is involved.\par\medskip
\begin{theorem}
\label{theorem-splitting}
Let $A$ be a finite central simple $k$-algebra.
Let $k'/k$ be a finite field extension.
The following are equivalent
\begin{enumerate}
\item $k'$ splits $A$, and
\item there exists a finite central simple algebra $B$ similar to $A$
such that $k' \subset B$ and $[B : k] = [k' : k]^2$.
\end{enumerate}
\end{theorem}

\begin{proof}
Assume (2). It suffices to show that $B \otimes_k k'$ is a matrix
algebra. We know that $B \otimes_k B^{op} \cong \text{End}_k(B)$.
Since $k'$ is the centralizer of $k'$ in $B^{op}$ by
Lemma \ref{lemma-self-centralizing-subfield}
we see that $B \otimes_k k'$ is the centralizer of $k \otimes k'$
in $B \otimes_k B^{op} = \text{End}_k(B)$. Of course this centralizer
is just $\text{End}_{k'}(B)$ where we view $B$ as a $k'$ vector space
via the embedding $k' \to B$. Thus the result.

\medskip\noindent
Assume (1). This means that we have an isomorphism
$A \otimes_k k' \cong \text{End}_{k'}(V)$ for some $k'$-vector space $V$.
Let $B$ be the commutant of $A$ in $\text{End}_k(V)$. Note that
$k'$ sits in $B$. By
Lemma \ref{lemma-when-tensor-is-equal}
the classes of $A$ and $B$ add up to zero in $\text{Br}(k)$.
From the dimension formula in
Theorem \ref{theorem-centralizer}
we see that
$$
[B : k] [A : k] =
\dim_k(V)^2 =
[k' : k]^2 \dim_{k'}(V)^2 =
[k' : k]^2 [A : k].
$$
Hence $[B : k] = [k' : k]^2$. Thus we have proved the result for the
opposite to the Brauer class of $A$. However, $k'$ splits the Brauer
class of $A$ if and only if it splits
the Brauer class of the opposite algebra, so we win anyway.
\end{proof}

\end{document}
